\section{Income tax accounting}

\begin{qfmethod}[Accounting profit, taxable income and temporary versus permanent differences]\label{meth:accounting-taxable-profit-differences}
\qffield{Statement} Accounting pretax profit follows reporting standards; taxable profit follows the applicable tax law. Current tax is the period's tax charge under that law, tax payable is the unsettled balance, and income-tax expense also includes deferred tax in profit or loss. A temporary difference arises from a difference between an asset/liability's carrying amount and tax base that reverses as it is recovered/settled. Permanent nondeductible expenses or exempt income affect the effective tax rate but do not create a reversing deferred tax balance \citep[IAS~12 paras.~7--8, 15, 24, 47 and~58]{adopted-accounting-standards-2023}.
\qffield{Practice} Separate illustrative tax case: accounting PBT 100 includes a nondeductible fine of 10 and warranty expense of 8; only 3 of warranties are paid/deductible this year. Tax depreciation exceeds book depreciation by 30. Taxable income is $100+10-30+8-3=85$. At 25 percent, current tax is 21.25. Faster tax depreciation creates DTL 7.5; the remaining warranty liability of 5 has tax base zero and creates DTA 1.25, assuming recovery is probable. Tax expense is $21.25+7.5-1.25=27.5$; net income 72.5. The fine explains the 2.5 excess over tax at the statutory rate, while the temporary differences explain cash timing.
\qffield{Interpretation} An asset with carrying amount 100 and tax base 70 generally has a taxable temporary difference of 30. A deductible accrued liability of 5 with tax base zero generally has a deductible difference of 5. Check recognition exceptions and the manner of recovery; do not memorize one sign rule without distinguishing assets from liabilities. The example assumes one jurisdiction, fixed rates, full recovery and no acquisitions or OCI.
\end{qfmethod}

\begin{qfmethod}[Deferred tax recognition, recovery and analytical treatment]\label{meth:deferred-tax-recovery-analysis}
\qffield{Statement} IFRS recognizes deductible-difference and loss-carryforward DTAs only to the extent taxable profit is probable, reviewing recovery each reporting date; US GAAP generally recognizes the DTA and applies a valuation allowance when realization is not more likely than not. Deferred taxes use expected reversal-period rates: IFRS permits enacted or substantively enacted rates, US GAAP enacted rates. They are not discounted in the reported accounts. Separate amounts recognized in profit, OCI, equity and acquisitions; total balance changes are not automatically deferred tax expense \citep[IAS~12 paras.~24, 34, 47, 53, 56, 58, 61A and~81]{adopted-accounting-standards-2023}; \citep[chap.~9, solutions~9--22, PDF pp.~207--208]{statement-analysis-workbook-2020}.
\qffield{Practice} The tax case's DTA 1.25 is recoverable only if deductions can be used in the right jurisdiction and before expiry. If only 60 percent is supportable, net recognized DTA becomes 0.75 and tax expense rises 0.50, without a current cash payment. A tax-rate reduction from 25 to 20 percent changes the depreciation DTL from 7.5 to 6 and the full warranty DTA from 1.25 to 1: net deferred-tax expense falls 1.25, assuming both originate in P\&L. A separate revaluation gain of 20 at 25 percent creates DTL 5 in OCI, not an extra 5 of profit tax or cash tax.
\qffield{Analysis} Read difference categories, expiration dates, valuation-allowance movements, rate changes and forecasts supporting losses. A release can lift earnings without improved operating cash flow. For valuation, forecast reversal cash taxes explicitly where material; deferred-tax liabilities are not automatically interest-bearing debt, nor automatically equity because they rolled forward historically. Growth may defer reversal, but disposal or lower reinvestment can bring it forward. Apply one cash-tax/valuation bridge and disclose sensitivity, rather than subtracting deferred tax twice.
\qffield{Source scope} SFAS 109 documents the US valuation-allowance and enacted-rate mechanism, now codified in ASC 740. This historical primary text is used with the 2020 workbook, not presented as the live consolidated Codification \citep[paras.~17--18 and~20--24, pp.~11--14]{us-deferred-tax-1992}.
\end{qfmethod}

\begin{qfmethod}[Statutory, effective and cash tax rates]\label{meth:tax-rate-reconciliation}
\qffield{Statement} The statutory rate is the applicable legal benchmark, the effective rate is reported income-tax expense divided by matching pretax income, and a stated cash tax rate is cash income tax paid divided by that pretax income. Some analysts use cash tax/taxable income instead; name the denominator. Near-zero or negative profit makes these ratios unstable. Geographic mix, permanent differences, credits, valuation allowances and discrete items explain effective-rate differences; temporary differences and settlements also explain cash differences \citep[IAS~12 paras.~58, 61A and~81]{adopted-accounting-standards-2023}.
\qffield{Practice} With PBT 100, current provision 21.25, net P\&L deferred expense 6.25 and total tax expense 27.5, effective rate is 27.5 percent versus statutory 25 percent. Opening tax payable 4 and closing payable 6 give cash tax $4+21.25-6=19.25$ and a 19.25 percent cash/PBT rate. It would be 22.65 percent against taxable income of 85, a different measure. The rate reconciliation is 25 percent plus 2.5 points from the nondeductible fine; the low cash rate is not evidence of a permanent tax advantage. An OCI-only DTL increase of 5 changes neither P\&L tax nor this payment ledger.
\qffield{Disclosure analysis} IAS 12 requires a tax-expense/rate reconciliation and temporary-difference information. US ASU 2023-09 expands annual rate reconciliations and cash-tax disaggregation, effective for public entities' annual periods beginning after 15 December 2024 and other entities after 15 December 2025. Its two 5 percent disaggregation thresholds use different bases: tax at the statutory rate for specified reconciliation items, and total net taxes paid for individual payment jurisdictions. Transition is prospective, with retrospective application permitted \citep[ASC~740-10-50-12A and~50-22--23; Summary p.~5]{us-tax-disclosures-2023}. Reconcile recurring jurisdictional taxes separately from nonrecurring releases before forecasting.
\end{qfmethod}
